\section{Conclusion}
\label{sec:conclusion}
\sys{} shows that per-request attribution of OS-induced delay does not
require a trap per request. Per-thread accounting kept by tracepoint
programs and read through a shared BPF map costs about as much as plain
event counters, and its accounting closes. The evaluation gives a more
nuanced answer than our hypothesis. Within a known workload, queueing
explains most tail requests, and any signal that tracks load (including
aggregate counts) detects them. Request-aligned kernel telemetry did
better on one held-out workload, but not when the held-out workload was
rotated or the runs repeated with fresh seeds, so its value lies in naming
the cause, not in detection. Controlled interventions moved
exactly the predicted delay component in seven of nine cases. The
remaining cases, together with the client-side receiver delay of the net
workload, show where the method stops: OS delay on threads outside the
request window, and host effects invisible to the guest. Future work is
to extend windows across all threads a request touches, to model
per-socket queues, and to repeat the study on bare metal. As a by-product,
we prepared a patch that lets bcc's \texttt{runqlat} separate wakeup
latency from preemption latency, the distinction our intervention results
relied on.
